% !TEX root = ../main.tex
\documentclass[../main.tex]{subfiles}

\begin{document}

%===============================================================================
% ENGLISH ABSTRACT
%===============================================================================
\chapter*{Abstract}
\addcontentsline{toc}{chapter}{Abstract}

This dissertation develops a multimodal cartography of role constellations in sustainability transitions, with empirical focus on intermediary organizations within the European Institute of Innovation and Technology (EIT) ecosystem. Drawing on Complex Adaptive Systems theory and Design Science Research methodology, we introduce \RoleBox---a computational toolkit for mapping how intermediaries position themselves and are positioned by others across multiple data modalities.

The research addresses a persistent methodological divide in transitions research: qualitative approaches offer rich insight but limited scale, while computational methods achieve scale but often miss theoretical depth. We propose multimodal cartography as a bridge---systematically integrating website texts, news media coverage, investment records, social media presence, and visual materials to triangulate role positions. These intermediaries connect to thousands of ecosystem actors; \RoleBox{} provides the infrastructure to scale such analysis across these networks.

Five empirical studies examine the role constellations of EIT intermediaries across these modalities. We find that interstitial positioning---intermediaries claiming to bridge separate communities---is substantially more common in self-presentation than in external positioning. Specifically, 34\% of intermediary actors claim connector roles, but only 12\% show cross-modal evidence for actually performing them. This ``interstitial gap'' reveals how role aspirations exceed role realities.

The dissertation contributes: (1) a theoretical framework for understanding role constellations as dynamic, multi-dimensional configurations; (2) \RoleBox, an open-source toolkit for multimodal role mapping that can scale to broader ecosystem analysis; (3) empirical findings on role dynamics among intermediaries in Europe's largest climate innovation ecosystem; and (4) methodological guidance for bridging qualitative and computational approaches in transitions research.

\vspace{0.5cm}

\noindent\textbf{Open Science Infrastructure}

All code, data pipelines, and interactive tools developed for this dissertation are publicly available:

\begin{itemize}[leftmargin=1.5cm, itemsep=0.2em]
\item \textbf{RoleBox} (umbrella repository): \texttt{github.com/babajideowoyele/rolebox}
\item \textbf{RoleSim} (agent-based simulation): \texttt{github.com/babajideowoyele/rolesim}
\item \textbf{RoleXray} (role extraction \& NLP): \texttt{github.com/babajideowoyele/rolexray}
\item \textbf{Data repositories}: Wikipedia/Wikidata, websites, Twitter, news, Crunchbase, images, video
\end{itemize}

\noindent Interactive data explorers with Carbon Design System UIs accompany each empirical chapter.

\vspace{0.3cm}

\textbf{Keywords:} role constellations, innovation ecosystems, sustainability transitions, multimodal analysis, complex adaptive systems, EIT, intermediaries, open science

\vspace{2cm}

%===============================================================================
% DUTCH ABSTRACT (Samenvatting)
%===============================================================================
\chapter*{Samenvatting}
\addcontentsline{toc}{chapter}{Samenvatting (Dutch Summary)}

Dit proefschrift ontwikkelt een multimodale cartografie van rolconstellaties in duurzaamheidstransities, met empirische focus op intermediaire organisaties binnen het ecosysteem van het Europees Instituut voor Innovatie en Technologie (EIT). Voortbouwend op de theorie van Complexe Adaptieve Systemen en Design Science Research methodologie, introduceren wij \RoleBox---een computationele toolkit voor het in kaart brengen van hoe intermediairs zichzelf positioneren en door anderen worden gepositioneerd over meerdere datamodaliteiten.

Het onderzoek adresseert een hardnekkige methodologische kloof in transitieonderzoek: kwalitatieve benaderingen bieden rijke inzichten maar beperkte schaal, terwijl computationele methoden schaal bereiken maar vaak theoretische diepgang missen. Wij stellen multimodale cartografie voor als brug---het systematisch integreren van websiteteksten, nieuwsmediaberichtgeving, investeringsgegevens, sociale media-aanwezigheid en visueel materiaal om rolposities te trianguleren. Deze intermediairs zijn verbonden met duizenden ecosysteemactoren; \RoleBox{} biedt de infrastructuur om dergelijke analyse over deze netwerken te schalen.

Vijf empirische studies onderzoeken de rolconstellaties van EIT-intermediairs over deze modaliteiten. Wij vinden dat interstitiële positionering---intermediairs die beweren afzonderlijke gemeenschappen te verbinden---substantieel vaker voorkomt in zelfpresentatie dan in externe positionering. Specifiek claimen 34\% van de intermediaire actoren verbindingsrollen, maar slechts 12\% toont cross-modale evidentie voor het daadwerkelijk uitvoeren ervan. Deze ``interstitiële kloof'' onthult hoe rolaspiraties rolrealiteiten overtreffen.

Het proefschrift draagt bij: (1) een theoretisch kader voor het begrijpen van rolconstellaties als dynamische, multidimensionale configuraties; (2) \RoleBox, een open-source toolkit voor multimodale rolkartering die kan schalen naar bredere ecosysteemanalyse; (3) empirische bevindingen over roldynamiek onder intermediairs in Europa's grootste klimaatinnovatie-ecosysteem; en (4) methodologische richtlijnen voor het overbruggen van kwalitatieve en computationele benaderingen in transitieonderzoek.

\textbf{Trefwoorden:} rolconstellaties, innovatie-ecosystemen, duurzaamheidstransities, multimodale analyse, complexe adaptieve systemen, EIT, intermediairs

\vspace{2cm}

%===============================================================================
% GERMAN ABSTRACT (Zusammenfassung)
%===============================================================================
\chapter*{Zusammenfassung}
\addcontentsline{toc}{chapter}{Zusammenfassung (German Summary)}

Diese Dissertation entwickelt eine multimodale Kartographie von Rollenkonstellationen in Nachhaltigkeitstransitionen, mit empirischem Fokus auf Intermediärorganisationen innerhalb des Ökosystems des Europäischen Instituts für Innovation und Technologie (EIT). Aufbauend auf der Theorie komplexer adaptiver Systeme und der Design Science Research Methodologie führen wir \RoleBox{} ein---ein computergestütztes Toolkit zur Erfassung, wie sich Intermediäre selbst positionieren und von anderen über mehrere Datenmodalitäten hinweg positioniert werden.

Die Forschung adressiert eine anhaltende methodische Kluft in der Transitionsforschung: Qualitative Ansätze bieten reichhaltige Einblicke, aber begrenzte Reichweite, während computergestützte Methoden Reichweite erzielen, aber oft theoretische Tiefe vermissen lassen. Wir schlagen multimodale Kartographie als Brücke vor---die systematische Integration von Website-Texten, Nachrichtenmedienberichterstattung, Investitionsdaten, Social-Media-Präsenz und visuellem Material zur Triangulation von Rollenpositionen. Diese Intermediäre sind mit Tausenden von Ökosystem-Akteuren verbunden; \RoleBox{} bietet die Infrastruktur, um solche Analysen über diese Netzwerke zu skalieren.

Fünf empirische Studien untersuchen die Rollenkonstellationen von EIT-Intermediären über diese Modalitäten hinweg. Wir stellen fest, dass interstitielle Positionierung---Intermediäre, die behaupten, getrennte Gemeinschaften zu verbinden---in der Selbstdarstellung wesentlich häufiger vorkommt als in der externen Positionierung. Konkret beanspruchen 34\% der Intermediär-Akteure Verbindungsrollen, aber nur 12\% zeigen modalitätsübergreifende Evidenz für deren tatsächliche Ausübung. Diese ``interstitielle Lücke'' zeigt, wie Rollenaspirationen die Rollenrealitäten übersteigen.

Die Dissertation leistet folgende Beiträge: (1) einen theoretischen Rahmen zum Verständnis von Rollenkonstellationen als dynamische, mehrdimensionale Konfigurationen; (2) \RoleBox, ein Open-Source-Toolkit für multimodale Rollenkartierung, das auf breitere Ökosystemanalysen skaliert werden kann; (3) empirische Befunde zur Rollendynamik unter Intermediären in Europas größtem Klimainnovations-Ökosystem; und (4) methodische Leitlinien zur Überbrückung qualitativer und computergestützter Ansätze in der Transitionsforschung.

\textbf{Schlüsselwörter:} Rollenkonstellationen, Innovationsökosysteme, Nachhaltigkeitstransitionen, multimodale Analyse, komplexe adaptive Systeme, EIT, Intermediäre

\end{document}
